\documentclass[11pt]{article}

% ---------- Core packages ----------
\usepackage[margin=1in]{geometry}
\usepackage{amsmath,amssymb}
\usepackage{graphicx}
\usepackage{booktabs}
\usepackage[numbers,sort&compress]{natbib}
\usepackage{hyperref}
\usepackage{caption}
\usepackage{authblk}
\usepackage{xcolor}

% ---------- Arabic bilingual abstract (requires XeLaTeX/LuaLaTeX + polyglossia) ----------
% Compile with: xelatex main.tex && bibtex main && xelatex main.tex && xelatex main.tex
\usepackage{fontspec}
\usepackage{polyglossia}
\setmainlanguage{english}
\setotherlanguage{arabic}
\newfontfamily\arabicfont[Script=Arabic]{Amiri}

\hypersetup{
  colorlinks=true,
  linkcolor=blue!50!black,
  citecolor=blue!50!black,
  urlcolor=blue!50!black
}

\title{\textbf{Physics-Guided Learning for Synthetic Temperature-Dependent Semiconductor Conductivity: Effects of Closed-Form Model Mismatch}}

\author[1]{Mohammed Hamzah Khudhur\thanks{Corresponding author: \texttt{mohammed.hamzah@uowasit.edu.iq}. ORCID: \href{https://orcid.org/0000-0002-7839-982X}{0000-0002-7839-982X}}}
\affil[1]{Department of Electrical Engineering, College of Engineering, University of Wasit, Kut, Iraq}
\date{}

\begin{document}
\maketitle

\begin{abstract}
\noindent
Accurate prediction of the temperature-dependent electrical conductivity $\sigma(T)$ of semiconductors is central to the design of power electronics, thermoelectric generators, and sensor materials, yet purely data-driven machine learning models routinely violate basic transport physics when extrapolated beyond their training temperature range. We introduce a physics-guided machine learning (PGML) architecture that couples electronic-structure descriptors -- band gap, carrier effective masses, density of states at the Fermi level, doping concentration, and Debye temperature -- with a differentiable transport layer built from standard relations of the Boltzmann transport equation under the relaxation-time approximation (Section~3.1 distinguishes what is rigorously derived from what is a phenomenological reduction of it): rather than mapping descriptors and temperature directly onto conductivity, the model first predicts a compact set of physically interpretable parameters governing carrier concentration and mobility separately, evaluates a closed-form transport expression from them, and applies a correction for higher-order effects. We test this architecture, a physics-regularized baseline (the identical physical-consistency loss applied to an unconstrained network, with no closed-form transport layer), a black-box multilayer perceptron, and two tree-ensemble baselines on a fully reproducible synthetic benchmark parameterized by electronic-structure descriptors of 24 real semiconductors (Si, Ge, GaAs, GaN, SiC, ZnO, InSb, PbTe, Bi$_2$Te$_3$, and related systems). We release the benchmark generator, all model and baseline implementations, training code, and per-seed results, and report seven successive revisions rather than only the best. A closed form outside the generator's functional family is the worst model on every metric, including physical consistency; revising it to match that family resolves consistency and closes most of the accuracy gap -- read cautiously, since this is closer to a demonstration of correct model specification than of general architectural benefit, given the deliberate functional-family match between the revised closed form and the benchmark's own generator. The remaining gap to Gradient-Boosted Trees survives hyperparameter tuning and is only partly closed by a physics-plus-boosting hybrid, which instead reaches the best extrapolation robustness of any single model tested on this synthetic benchmark. Combining every model tried, weighted by validation-set accuracy, reaches numerically lower error and better extrapolation than tuned Gradient-Boosted Trees, but neither an exploratory paired test across the 5 seeds ($p > 0.1$, paired $t$-test) nor an exploratory material-level bootstrap over the benchmark's materials supports this gap as statistically significant on any metric; since the combination includes Gradient-Boosted Trees itself as a member, this is in any case not evidence that the physics-guided architecture outperforms it independently. We argue that architectural-physics-benefit claims should be checked against multiple closed-form complexities, both black-box and physics-regularized baselines, and paired significance tests across seeds before being generalized -- and that none of these synthetic-benchmark findings should be read as claims about real, measured semiconductor conductivity data.

\vspace{0.5em}
\noindent\textbf{Keywords:} physics-informed machine learning, semiconductor conductivity, electronic-structure descriptors, Boltzmann transport equation, materials informatics, temperature-dependent transport properties
\end{abstract}

\begin{otherlanguage}{arabic}
\begin{abstract}
\noindent
غالبًا ما يُفترض أن دمج الفيزياء معماريًا في نماذج التعلّم الآلي يحسّن التنبؤ بالخصائص، دون اختبار هذا الافتراض عبر أكثر من صيغة مغلقة واحدة، أو مقابل مرجعين مختلفين معًا: نموذج "صندوق أسود" ونموذج مُنظَّم فيزيائيًا فقط عبر دالة الخسارة. نبني معيارًا قياسيًا اصطناعيًا قابلًا للتكرار الكامل: منحنيات توصيلية كهربائية معتمدة على درجة الحرارة $\sigma(T)$، مُولَّدة من نموذج ظاهري مبسّط \emph{مستوحى} من معادلة بولتزمان للنقل -- وليس اشتقاقًا صارمًا منها -- ومُعامَل بواصفات إلكترونية بنيوية لِـ24 مادة شبه موصلة حقيقية، علمًا أن قيم $\sigma(T)$ نفسها اصطناعية بالكامل وليست قياسات فعلية. نختبر على هذا المعيار بِنيةً موجَّهة فيزيائيًا (PGML) تتنبأ بمعاملات نقل قابلة للتفسير عبر طبقة نقل تحليلية مغلقة، مقابل شبكة عصبية "صندوق أسود"، ونموذج مُنظَّم فيزيائيًا (دالة الخسارة الفيزيائية نفسها دون طبقة مغلقة)، ونموذجَي أشجار قرار تجميعيَّين، مع تقسيم صارم على مستوى المادة (لا تظهر أي مادة في أكثر من مجموعة واحدة) وتقسيم إضافي يستبعد درجات الحرارة العالية أثناء التدريب لاختبار التعميم على مواد ودرجات حرارة غير مرئية معًا. نعرض سبع تجارب متعاقبة بدلًا من الاقتصار على الأفضل أداءً. صيغة مغلقة خارج العائلة الدالية لمولّد البيانات كانت الأضعف على كل مقياس، بما فيها الاتساق الفيزيائي؛ وإعادة صياغتها لتطابق تلك العائلة تحل مشكلة الاتساق وتُغلق معظم فجوة الدقة -- وهذه نتيجة ينبغي قراءتها بحذر، إذ هي أقرب إلى إثبات مطابقة صحيحة للنموذج منها إلى دليل عام على فائدة الدمج المعماري، نظرًا للتطابق المتعمد بين الصيغة المُعاد صياغتها ومولّد المعيار القياسي نفسه. الفجوة المتبقية مع أشجار التعزيز التدرجي تصمد أمام ضبط المعاملات، ولا يُغلقها إلا جزئيًا مزيج فيزيائي-تعزيزي يحقق بالمقابل أفضل متانة تعميم بين كل النماذج المفردة. ودمج كل نموذج جُرِّب، مُرجَّحًا بدقته على مجموعة التحقق، يحقق خطأ أقل وتعميمًا أفضل رقميًا من أشجار التعزيز التدرجي المضبوطة، لكن الاختبارات الإحصائية المزدوجة عبر البذور الخمس لا تدعم هذا الفارق كدلالة إحصائية على أي مقياس ($p > 0.1$، اختبار $t$ مزدوج)؛ ولأن هذا المزيج يضم أشجار التعزيز التدرجي نفسها كعضو فيه، فهو بأي حال ليس دليلًا على تفوّق البِنية الفيزيائية بمفردها. نتيح مولّد المعيار القياسي وكل تطبيقات النماذج والأكواد ونتائج كل بذرة على حدة، ونرى أن ادّعاءات فائدة الدمج المعماري للفيزياء ينبغي أن تُختبر مقابل أكثر من صيغة مغلقة، ومقابل النوعين المرجعيين معًا، وباختبارات دلالة إحصائية مزدوجة عبر البذور، قبل تعميمها -- وأن لا شيء من نتائج هذا المعيار الاصطناعي ينبغي أن يُقرأ كادّعاء بشأن بيانات ناقلية أشباه الموصلات الحقيقية المقيسة.

\vspace{0.5em}
\noindent\textbf{الكلمات المفتاحية:} التعلّم الآلي الموجَّه فيزيائيًا، توصيلية أشباه الموصلات، الواصفات الإلكترونية البنيوية، معادلة بولتزمان للنقل، معلوماتية المواد، خصائص النقل المعتمدة على درجة الحرارة
\end{abstract}
\end{otherlanguage}

\section{Introduction}

Electrical conductivity $\sigma(T)$ governs the operating envelope of nearly every semiconductor device, from power diodes and high-temperature sensors to thermoelectric generators that convert waste heat into electricity. Because $\sigma(T)$ results from the interplay of carrier concentration, scattering mechanisms, and band structure, it is highly nonlinear and material-specific: intrinsic semiconductors show conductivity rising exponentially with temperature as thermally generated carriers cross the band gap, while heavily doped or degenerate semiconductors show conductivity falling with temperature as phonon scattering intensifies, mirroring metallic behavior. Predicting $\sigma(T)$ accurately across this range of regimes, and for materials that have not yet been synthesized or measured, would substantially accelerate the design of power electronics and thermoelectric modules \citep{gorai2017computationally,snyder2008complex}.

Two complementary strategies exist for this prediction task. The first-principles route solves or parameterizes the Boltzmann transport equation (BTE) using band structures and scattering rates obtained from density functional theory (DFT), yielding physically consistent predictions but at a computational cost that scales poorly with the size of the chemical space to be screened. The data-driven route trains a regression model -- typically a neural network, random forest, or gradient-boosted tree ensemble -- directly on measured or computed conductivity data, using materials descriptors as features \citep{ward2016general,butler2018machine,ramprasad2017machine}. This approach is fast and scales to large candidate libraries, but because it has no knowledge of the governing transport equations, it can produce predictions that are numerically plausible yet physically absurd: conductivity that decreases with temperature in the intrinsic regime, or that fails to follow the expected activation-energy slope on an Arrhenius plot, especially when extrapolating beyond the temperature range seen during training.

Physics-informed machine learning (PIML) has emerged as a middle path, injecting governing-equation knowledge into the training process rather than relying on data alone \citep{raissi2019pinn,karniadakis2021pgml,willard2020integrating}. Most PIML work to date targets continuous field problems -- fluid flow, heat conduction, wave propagation -- where a partial differential equation residual can be evaluated directly on network outputs via automatic differentiation. Materials property prediction poses a different structure: the underlying physics (the BTE under the relaxation-time approximation) is well established in reduced, closed-form limits, but the microscopic parameters that enter it (scattering exponents, effective activation energies, prefactor mobilities) are not directly observable and must themselves be inferred from electronic-structure descriptors. This motivates a design choice, examined in this paper, in which the machine learning model does not predict conductivity directly, and does not merely add a soft physics penalty to an otherwise unconstrained network, but instead predicts the physical transport parameters themselves and evaluates conductivity through the transport equation as an explicit computational layer.

We refer to this design as physics-guided machine learning (PGML) to distinguish it from the more common physics-informed neural network (PINN) pattern of penalizing PDE residuals in the loss function of an otherwise black-box network. In principle, a transport equation evaluated as a forward layer is satisfied by construction for any values the physics-parameter head produces, whereas a soft penalty can be traded off against data fit during optimization and satisfied only approximately -- an architectural advantage frequently assumed in the physics-informed machine learning literature. Whether that advantage actually materializes for a given reduced physical law and a given data-generating process is an empirical question, and testing it honestly requires comparing against both a black-box baseline and a loss-only physics-informed baseline on the same data. This paper makes nine contributions.

\begin{enumerate}
\item We formulate a physics-guided architecture in which a descriptor-and-temperature encoder predicts interpretable transport parameters governing carrier concentration and mobility separately, which are passed through a closed-form, BTE-inspired transport expression for $\sigma(T)$ built from standard BTE-derived governing relations, followed by a correction for higher-order effects.
\item We define a composite training objective combining log-conductivity data loss, a shape-consistency penalty on the sign and curvature of $d\sigma/dT$ per regime, and a boundedness penalty that keeps any residual correction small relative to the physical baseline.
\item We construct and publicly release a fully reproducible synthetic benchmark of 24 real semiconductor systems, together with the complete model, baseline, and training code, and use it to compare the proposed architecture against black-box, tree-ensemble, and loss-only physics-informed (residual-only PRN) baselines under identical descriptors and evaluation protocol.
\item We show, by comparing closed forms of differing functional richness on this benchmark, that the benefit of architectural physics integration is highly sensitive to how closely the closed form's functional family matches the data-generating process: a more reduced closed form is the worst-performing model on every metric, including the physical-consistency violation rate it is designed to improve, while a closed form that keeps the transport equation's carrier-concentration and mobility terms separate resolves the consistency problem entirely and closes most of the accuracy gap to the strongest baselines -- a methodological lesson for how architectural-physics-integration claims should be evaluated.
\item We further show that replacing a generic bounded residual-correction network with an explicit second conduction channel, structurally matched to the benchmark's own narrow-gap high-temperature correction rather than left to an unconstrained network of arbitrary shape, closes nearly all of the remaining accuracy gap to the strongest neural baseline -- evidence that committing to the \emph{right} functional form, not just adding architectural flexibility, is what drives the benefit.
\item We test whether the remaining gap to tree-ensemble baselines is closable at all, three ways: tuning Gradient-Boosted Trees' own hyperparameters via a validation-selected grid search barely moves its accuracy, ruling out an under-tuned baseline as the explanation; fitting Gradient-Boosted Trees on the closed form's residual narrows the in-distribution accuracy gap without closing it, but achieves, on this synthetic benchmark, the best extrapolation performance of every model in the paper; and Gradient-Boosted Trees given the closed form's fitted parameters as extra input features, instead of its residual, does not help at all, isolating the closed form's output prediction, not its internal parameters, as what makes the residual hybrid work, and suggesting the closed form's practical advantage on this benchmark is extrapolation robustness rather than in-distribution accuracy.
\item We show that integrating every model tried into one predictor -- a simple unweighted average of all eight base models -- comes numerically closer to closing the gap than any single model or hybrid, with a mean at or ahead of the strongest individual baseline on both accuracy and extrapolation (though not distinguishably so at $n=5$ seeds), while a stacked variant that learns combination weights from the (small) validation split does not improve on it, because the learned weights are unstable across seeds -- evidence that on a benchmark this small, integration without additional data-driven weight estimation outperforms integration with it.
\item We test whether that ensemble was leaving gains on the table by excluding the two physics-plus-boosting hybrids as candidate members: extending both combination strategies to all ten models leaves the simple average statistically unchanged and gives the stacked variant only a small, still-insufficient improvement, showing that the bottleneck is the combination strategy and the size of the validation split available to fit it, not which models are eligible to be combined.
\item We show that a combination strategy between the two extremes already tried -- weighting each base model's prediction by the inverse of its own squared validation RMSE, rather than averaging uniformly or fitting a fully free regression -- has a lower mean error than both on every metric, and is the only result in the paper with a mean simultaneously more accurate and more extrapolation-robust than tuned Gradient-Boosted Trees while matching its $R^2$ within noise, because it needs far less from the same small validation split than a joint regression over correlated candidates does; neither an exploratory paired test across the 5 seeds nor an exploratory material-level bootstrap supports any of these differences as statistically significant, and the comparison includes Gradient-Boosted Trees itself as an ensemble member.
\end{enumerate}

The remainder of the paper is organized as follows. Section~2 reviews related work in materials informatics and physics-informed learning. Section~3 details the physics background, descriptor set, model architecture, and loss function. Section~4 describes the benchmark and experimental protocol. Section~5 presents results and ablations. Section~6 discusses implications and failure modes, Section~7 states limitations, and Section~8 concludes.

\section{Related Work}

\subsection{Materials Informatics and Descriptor-Based Property Prediction}

Materials informatics has established descriptor-based machine learning as a practical alternative to exhaustive first-principles screening. Ward et al. introduced a general-purpose feature set (composition-, structure-, and electronic-structure-derived descriptors) for predicting formation energies and other properties across broad chemical spaces \citep{ward2016general}. Large open databases such as the Materials Project \citep{jain2013materialsproject} and AFLOW \citep{curtarolo2012aflow} provide the DFT-computed band structures, densities of states, and effective masses that make such descriptors available at scale, and have been the substrate for numerous property-prediction studies reviewed by Butler et al. and Ramprasad et al. \citep{butler2018machine,ramprasad2017machine}. Graph-based representations that learn directly from crystal structure, such as CGCNN \citep{xie2018crystal} and MEGNet \citep{chen2019graph}, and multi-fidelity extensions that fuse DFT and experimental labels \citep{chen2021learning}, have pushed accuracy higher for properties such as formation energy and band gap. Descriptor-based machine learning has also been applied directly to transport-adjacent properties: Miyazaki et al. predict half-Heusler compounds' lattice thermal conductivity from atomic-radius and atomic-mass descriptors of the constituent elements alone \citep{miyazaki2021thermal}, and Antunes et al. use an attention-based deep network over compositional descriptors to predict room-temperature thermoelectric transport coefficients (Seebeck coefficient, electrical conductivity, thermal conductivity) directly from composition \citep{antunes2023thermoelectric}. Both, like the present work, target properties governed by the same phonon- and carrier-scattering physics as electrical conductivity, but predict a single value (typically at or near 300~K) rather than a full temperature-dependent curve spanning the intrinsic-to-extrinsic crossover; extending descriptor-based ML to the explicit $T$-dependence of $\sigma$, rather than a room-temperature snapshot, is one of the gaps the present benchmark targets. Zhang and Ling highlight the small-data regime typical of experimentally measured transport properties and argue for inductive biases that reduce the effective hypothesis space \citep{zhang2018strategy}, a concern directly relevant to conductivity prediction, where high-quality temperature-resolved measurements are comparatively scarce relative to, e.g., computed formation energies.

\subsection{Physics-Informed and Physics-Guided Machine Learning}

Physics-informed neural networks embed differential-equation residuals into the training loss so that a network trained on sparse data remains consistent with known governing equations \citep{raissi2019pinn}. Karniadakis et al. survey the broader physics-informed machine learning landscape, distinguishing approaches that inject physics through the loss function, through the network architecture, through data augmentation, or through hybrid model structures \citep{karniadakis2021pgml}. Willard et al. similarly catalog strategies for integrating physics-based and data-driven modeling, noting that architectural integration -- where physical structure constrains the hypothesis space directly rather than only through a penalty term -- tends to generalize better under distribution shift than loss-only approaches \citep{willard2020integrating}. Closest to the present work's application domain, Pimachev and Neogi train an electronic-transport-informatics framework on \emph{ab initio} simulations of small semiconductor systems and use it to predict the thermopower of fabricated silicon/germanium heterostructures too large to simulate directly \citep{pimachev2021transport}; their framework is physics-\emph{aware} in the sense of being trained on physics-based simulation outputs rather than a differentiable closed-form physics layer, which is the specific architectural-integration mechanism this paper isolates and tests. Our design follows the architectural-integration strategy Willard et al.\ describe: the transport equation is not merely a training-time penalty but a forward computational layer that the network's outputs must pass through, which is closer in spirit to physics-guided or physics-constrained architectures than to a canonical PINN.

\subsection{Transport Theory for Semiconductor Conductivity}

The relaxation-time approximation to the Boltzmann transport equation remains the standard reduced description of semiconductor conductivity, expressing $\sigma$ as the product of carrier concentration, relaxation time, and inverse effective mass \citep{ashcroftmermin1976}. Deformation-potential theory quantifies acoustic-phonon-limited mobility and its characteristic $T^{-3/2}$ temperature dependence \citep{bardeen1950deformation}, while ionized-impurity scattering produces the complementary $T^{3/2}$ dependence that dominates at low temperature in doped materials. Thermoelectric materials research has long used these reduced transport models, together with band-structure descriptors, to screen candidate compounds computationally \citep{gorai2017computationally,snyder2008complex}, but typically as forward physics models rather than as a differentiable layer inside a trainable machine learning system. The present work adapts this transport theory into exactly such a differentiable layer.

\section{Methodology}

\subsection{Physical Background}

Under the relaxation-time approximation, the electrical conductivity of a semiconductor can be written as
\begin{equation}
\sigma(T) = n(T)\, e\, \mu(T) = n(T)\, e^{2}\, \tau(T) / m^{*},
\label{eq:bte}
\end{equation}
where $n(T)$ is the carrier concentration, $e$ the electron charge, $\mu(T)$ the mobility, $\tau(T)$ the momentum relaxation time, and $m^{*}$ the conductivity effective mass. In the intrinsic regime, carrier concentration follows
\begin{equation}
n_i(T) = 2\left(\frac{2\pi k_B T}{h^2}\right)^{3/2} \left(m_e^{*} m_h^{*}\right)^{3/4} \exp\!\left(-\frac{E_g}{2 k_B T}\right),
\label{eq:ni}
\end{equation}
which produces an Arrhenius-type exponential rise of $\sigma$ with $T$, with a characteristic slope of $-E_g/2$ on a $\ln\sigma$ vs.\ $1/T$ plot. In the extrinsic (doped) regime, carrier concentration saturates near the dopant density and conductivity is instead governed by the temperature dependence of $\tau(T)$, combining acoustic-phonon scattering ($\tau_{\text{ac}} \propto T^{-3/2}$) \citep{bardeen1950deformation} and ionized-impurity scattering ($\tau_{\text{ii}} \propto T^{3/2}$) via Matthiessen's rule, $1/\tau = 1/\tau_{\text{ac}} + 1/\tau_{\text{ii}}$. Equations~\ref{eq:bte}--\ref{eq:ni} and the scattering exponents above are standard, textbook results under the relaxation-time approximation. The closed-form transport layer we build from them below (Equations~\ref{eq:n_closed}--\ref{eq:sigma_phys}) instead fits its material-specific prefactors and effective parameters to data rather than computing them from microscopic theory for each material; we accordingly describe the overall closed form as BTE-\emph{inspired} rather than BTE-\emph{derived}, reserving the latter term for the governing relations themselves.

A single combined exponential-times-power-law closed form -- $\sigma_{\text{phys}} \propto n_0 T^{3/2}\exp(-E_{g,\text{eff}}/2k_BT)\,\tau_0(T/T_0)^{-p}$, as in an earlier version of this work -- commits to one effective activation energy and one effective scattering exponent per material, which cannot separately represent the extrinsic plateau and the intrinsic rise, nor two scattering channels with different temperature exponents; we found empirically that this cost more accuracy, extrapolation reliability, and physical consistency than it gained (Section~5.1). We instead keep Equations~\ref{eq:bte}--\ref{eq:ni}'s two physical ingredients -- carrier concentration and mobility -- as separate closed-form pieces, each with its own free parameters, so the closed form is in the same functional family as the transport process it approximates rather than a further-reduced simplification of it. Carrier concentration uses the non-degenerate charge-neutrality approximation directly:
\begin{equation}
n(T;\theta) = \frac{N_0}{2} + \sqrt{\left(\frac{N_0}{2}\right)^2 + n_i(T;\theta)^2}, \qquad
n_i(T;\theta) = C_n \left(\frac{T}{T_0}\right)^{3/2} \exp\!\left(-\frac{E_{g,\text{eff}}}{2 k_B T}\right),
\label{eq:n_closed}
\end{equation}
which smoothly interpolates between an extrinsic floor $N_0$ (the fitted, doping-set carrier density) and the intrinsic exponential rise, exactly as in Equation~\ref{eq:ni} but with $C_n$ and $E_{g,\text{eff}}$ free rather than fixed to their DFT values. Mobility keeps the two scattering channels of Equation~\ref{eq:bte} separate and combines them via Matthiessen's rule:
\begin{equation}
\frac{1}{\mu(T;\theta)} = \frac{1}{\tau_{\text{ac}0}}\left(\frac{T}{T_0}\right)^{3/2} + \frac{1}{\tau_{\text{ii}0}}\left(\frac{T}{T_0}\right)^{-3/2},
\label{eq:mu_closed}
\end{equation}
with independent learned prefactors $\tau_{\text{ac}0}$ (acoustic-phonon channel) and $\tau_{\text{ii}0}$ (ionized-impurity channel) rather than one combined exponent. The closed-form transport expression is then
\begin{equation}
\sigma_{\text{phys}}(T;\theta) = n(T;\theta)\,\mu(T;\theta) / m^{*}, \qquad \theta = (E_{g,\text{eff}}, C_n, N_0, \tau_{\text{ac}0}, \tau_{\text{ii}0}),
\label{eq:sigma_phys}
\end{equation}
five physically interpretable parameters per material in place of the earlier four, none of which are fixed to their nominal descriptor-derived values -- $E_{g,\text{eff}}$ and $C_n$ absorb any doping-induced shift and any DFT-to-experiment gap correction, and $N_0$, $\tau_{\text{ac}0}$, $\tau_{\text{ii}0}$ are fitted rather than read off a specific microscopic scattering theory.

\subsection{Electronic-Structure Descriptors}

For each material and doping condition, we assemble a descriptor vector $x \in \mathbb{R}^{9}$:
\begin{itemize}
\item Band gap $E_g$ (eV): a representative literature-typical value for each material, compiled from standard semiconductor-physics references rather than computed via DFT for this work; not claimed exact for any single sample or source.
\item Electron and hole effective masses $m_e^{*}$, $m_h^{*}$ (units of free-electron mass $m_0$): likewise representative literature-typical values, simplified single-valley conductivity-relevant averages rather than precise multi-valley density-of-states masses computed here.
\item A density-of-states-at-the-Fermi-level \emph{proxy}, $\text{DOS}(E_F)_{\text{proxy}} = \sqrt{m_e^{*} m_h^{*}} \cdot \log_{10} n(300\,\text{K}; N_d)$ (dimensionless engineered feature, not a physical density of states in states/eV/atom): grows with effective mass and with the 300~K free-carrier density implied by the doping level $N_d$, intended only as an ML input feature, not a DFT-computed quantity.
\item Nominal carrier (doping) concentration $N_d$ ($\text{cm}^{-3}$), on a log scale: one of three discrete values per material (intrinsic, $10^{17}$, $10^{19}\,\text{cm}^{-3}$), held fixed rather than varying continuously or itself uncertain.
\item Debye temperature $\Theta_D$ (K): a representative literature-typical value, used as a proxy for phonon scattering strength.
\item Average atomic mass $\bar{M}$ and mean electronegativity difference $\Delta\chi$ of the constituent elements, as compact proxies for deformation-potential and polar-optical scattering strength; unlike the other descriptors above, these two are computed exactly from each element's standard Pauling electronegativity and atomic weight, not looked up per material.
\item Lattice thermal conductivity $\kappa_L$ (W/m/K) at 300~K: a representative literature-typical value, included because materials with strong anharmonic phonon scattering (relevant to $\tau_{\text{ac}}$) also tend to show suppressed $\kappa_L$.
\end{itemize}
Except for $\bar{M}$ and $\Delta\chi$, these descriptors are deliberately simplified, literature-typical inputs suited to a reproducible demonstration benchmark, not original per-material electronic-structure calculations; Section~\ref{sec:limitations} returns to what this does and does not license us to conclude.
Temperature $T$ is supplied separately from the static descriptors and is encoded both as a raw value (min--max normalized over 80--800~K) and as $1/T$, since transport behavior is naturally expressed in Arrhenius coordinates.

\subsection{Physics-Guided Architecture}

The model consists of the physics-parameter head and closed-form transport layer described below, plus one of two mechanisms -- an explicit second conduction channel or a generic residual correction network, ablated against each other in Section~5.3 -- for handling transport physics the base five-parameter closed form does not represent.

\textbf{Physics-Parameter Head (PPH).} A four-layer multilayer perceptron with widths $(128, 128, 64, n_\theta)$ and GELU activations maps the descriptor vector $x$ (temperature-independent) to raw outputs, which are then transformed by output activations that guarantee physical admissibility: a $\text{softplus}(\cdot) + \epsilon$ activation for each of the five base parameters in Equation~\ref{eq:sigma_phys} (non-negative by construction), and, when the second conduction channel is enabled ($n_\theta = 8$ rather than $5$), a $\text{sigmoid}(\cdot)$ activation bounding the channel's amplitude to $[0,1]$ and a scaled sigmoid placing its turn-on temperature in $[50, 950]\,\text{K}$ (Section~3.3, second channel, below). Because the PPH does not see $T$ directly, its outputs are guaranteed to be temperature-independent material constants, matching the physical role of these quantities in Equations~\ref{eq:n_closed}--\ref{eq:sigma_phys}.

\textbf{Closed-Form Transport Layer.} Given $\theta = (E_{g,\text{eff}}, C_n, N_0, \tau_{\text{ac}0}, \tau_{\text{ii}0})$ from the PPH and a query temperature $T$, this parameter-free layer evaluates Equation~\ref{eq:sigma_phys} exactly, using $k_B$, $T_0 = 300$~K, and $m^{*} = m_e^{*} + m_h^{*}$ as fixed constants. Because this layer contains no trainable weights, gradients with respect to $\theta$ flow directly through a known analytic function, which keeps the physics constraint exact rather than approximate.

\textbf{Explicit Second Conduction Channel.} Section~3.5's benchmark generator applies a smooth high-temperature multiplicative correction to narrow-gap materials, $1 + \gamma\,\text{sigmoid}((T-T_c)/T_w)$, representing a second conduction channel (e.g.\ a second band or valley) that the base closed form's single carrier-concentration and mobility terms do not capture. Rather than leave this to a generic correction network, we give the PPH three additional free parameters -- amplitude $\gamma_2 \in [0,1]$, turn-on center $T_{c2}$, and turn-on width $T_{w2}$ -- and apply the structurally identical correction to the closed-form output:
\begin{equation}
\sigma_{\text{pred}}(T;\theta) = \sigma_{\text{phys}}(T;\theta)\,\big(1 + \gamma_2\,\text{sigmoid}((T-T_{c2})/T_{w2})\big).
\label{eq:second_channel}
\end{equation}
Unlike the generator, $\gamma_2$, $T_{c2}$, and $T_{w2}$ are learned per material from descriptors rather than fixed to $E_g < 0.5\,\text{eV}$; a material with no such channel is expected to converge toward $\gamma_2 \approx 0$. This differs from a generic residual correction only in that it commits to the generator's actual functional form (a single logistic turn-on) rather than an unconstrained correction of arbitrary shape.

\textbf{Residual Correction Network (RCN).} As an alternative -- or, in one ablation arm, an addition -- to the explicit second channel, a second, smaller network (widths $(64,32,1)$) takes as input the concatenation of $x$, $T$, and $\log \sigma_{\text{phys}}(T;\theta)$ (or $\log \sigma_{\text{pred}}(T;\theta)$ from Equation~\ref{eq:second_channel} when both mechanisms are enabled), and outputs a scalar $\delta \in [-1,1]$ via a $\tanh$ output activation. The final prediction is
\begin{equation}
\sigma_{\text{pred}}(T) = \sigma_{\text{phys or channel-corrected}}(T;\theta)\, \big(1 + \eta\, \delta(x,T)\big),
\label{eq:final}
\end{equation}
with $\eta = 0.15$ fixed as a hyperparameter bounding the maximum multiplicative correction to $\pm 15\%$. Unlike the second channel, the RCN's correction has no particular functional shape and can in principle absorb any effect the closed form omits -- at the cost of the interpretability and inductive bias that come with committing to a specific functional form.

\subsection{Physics-Guided Loss Function}

Training minimizes a composite objective over a batch of $(x_i, T_{ij}, \sigma_{ij})$ triples, where $j$ indexes multiple temperature points per material-doping condition $i$:
\begin{equation}
\mathcal{L} = \mathcal{L}_{\text{data}} + \lambda_1 \mathcal{L}_{\text{shape}} + \lambda_2 \mathcal{L}_{\text{bound}}.
\end{equation}

\textbf{Data term.} $\mathcal{L}_{\text{data}} = \frac{1}{N}\sum_{i,j} \big(\log_{10}\sigma_{\text{pred}}(T_{ij}) - \log_{10}\sigma_{ij}\big)^2$, fit in log space because $\sigma$ spans many orders of magnitude across the intrinsic-to-degenerate range covered by the benchmark.

\textbf{Shape-consistency term.} For each curve $i$, we numerically differentiate the model's predicted $\log\sigma$ with respect to $1/T$ at the sampled points (via a finite-difference gradient over the shared temperature grid) and convert the local slope to an implied activation energy, $E_{a,\text{pred}} = -2k_B \times (\text{slope})$, which is $O(1)\,\text{eV}$ and directly comparable to the descriptor's own band gap -- working in raw slope units instead, which are $O(1/2k_B)\sim 6{,}000\,\text{K}$, made the squared-error term dominate the total loss by orders of magnitude in early experiments. In temperature windows classified (from the labeled dopant density and the descriptor's own $E_g$, $m_e^{*}$, $m_h^{*}$) as intrinsic-dominated ($n_i(T) > 3 N_d$) and measurable (Section~3.5), $\mathcal{L}_{\text{shape}}$'s primary term is the mean squared deviation between $E_{a,\text{pred}}$ and $E_g$. A secondary term penalizes $|d\log\sigma/d\log T|$ exceeding a generic bound of 3 anywhere on the curve, discouraging unphysically steep local slopes without committing to a specific extrinsic-regime target.

\textbf{Boundedness term.} $\mathcal{L}_{\text{bound}} = \frac{1}{N}\sum_{i,j} \delta(x_i, T_{ij})^2$, an $L_2$ penalty on the residual-correction output that discourages the RCN from doing more work than necessary, keeping $\sigma_{\text{phys}}$ as the dominant contributor whenever the data do not require otherwise.

\textbf{Regularization term.} $\mathcal{L}_{\text{reg}}$ is standard weight decay applied to all trainable parameters via the optimizer (AdamW, $10^{-4}$), rather than a separate additive loss term; we accordingly write the composite objective without an explicit $\lambda_3\mathcal{L}_{\text{reg}}$ term below.

We set $\lambda_1 = 0.02$ and $\lambda_2 = 0.05$ from a small manual sweep of $\lambda_1 \in \{0, 0.02, 0.05, 0.1\}$ on one seed's validation split with the richer closed form (Section~5.2); $\lambda_1=0.02$ gave the lowest validation loss; a larger $\lambda_1=0.1$, appropriate for the more reduced closed form used in an earlier iteration (Section~5.1), degraded validation accuracy once the richer form was adopted, illustrating that the right amount of loss-level physics guidance depends on how much the architecture already gets right on its own.

\subsection{Benchmark Construction}

Rather than digitizing individual literature measurements, which cannot be redistributed at scale and are difficult to verify sample-by-sample, we construct a fully reproducible synthetic benchmark and release the generator, the model implementation, and every experiment script alongside this paper.\footnote{Code: \texttt{papers/35-physics-guided-ml-semiconductor-conductivity/code/} in the accompanying repository.} We compile electronic-structure descriptors for 24 real semiconductors spanning elemental, III-V, II-VI, IV-IV, and narrow-gap chalcogenide/telluride systems (Si, Ge, diamond, GaAs, GaN, GaP, GaSb, InP, InAs, InSb, AlAs, AlSb, AlN, ZnO, ZnS, ZnSe, CdS, CdSe, CdTe, PbTe, PbSe, PbS, 4H-SiC, Bi$_2$Te$_3$), using representative literature values for band gap, conductivity effective masses, Debye temperature, and lattice thermal conductivity, and computing electronegativity difference and average atomic mass exactly from the constituent elements' Pauling electronegativities and atomic weights. Each material is instantiated at three nominal doping levels (intrinsic, $N_d=10^{17}\,\text{cm}^{-3}$, $N_d=10^{19}\,\text{cm}^{-3}$), giving 72 material-doping curves.

For each curve, $\sigma(T)$ is \emph{synthesized} at 25 temperatures over 80--800~K from a reduced transport model that shares the functional family of Equations~\ref{eq:n_closed}--\ref{eq:mu_closed} -- the same charge-neutrality carrier-concentration blend and the same two Matthiessen-combined scattering channels -- but with its own material-specific parameter values ($N_d$ from the nominal doping level rather than a fitted $N_0$; $C_n$, $\tau_{\text{ac}0}$, $\tau_{\text{ii}0}$ tied to the descriptor table's band gap, effective masses, and Debye temperature via the fixed physical relations in Section~3.1, rather than fitted freely). Narrow-gap materials ($E_g < 0.5\,\text{eV}$: InAs, InSb, PbTe, PbSe, PbS, Bi$_2$Te$_3$) additionally receive a smooth high-temperature multiplicative correction (a logistic turn-on above 500~K, $\pm 20\%$), representing transport physics -- e.g.\ a second conduction channel -- that lies outside \emph{any} two-channel Matthiessen closed form, model or generator alike. This gives the physics-parameter head a data-generating process it can in principle match almost exactly (Section~6 confirms it does, up to this narrow-gap correction and the noise below), while still leaving a genuine test of whether the residual correction network and the shape-consistency loss help close that remaining, deliberately unmatched gap. Each curve additionally receives a per-curve log-normal scale factor ($\sigma_{\ln}=0.08$) and per-point log-normal noise ($\sigma_{\ln}=0.05$), representing sample-to-sample variability and measurement noise. Points below $10^{-9}\,\text{S/cm}$ -- below which a sample reads as immeasurably insulating on typical characterization equipment -- are excluded from training and evaluation rather than fit at an arbitrary floor value, since a hard floor would introduce a non-physical kink that no smooth transport law can represent; after this filter the benchmark contains 1{,}618 (material, doping, temperature, conductivity) samples.

We split at the \emph{material} level (70\%/15\%/15\% of the 24 materials, all three doping conditions of a given material kept together) so that test-time evaluation always requires generalizing to materials never seen during training, not merely to a new doping level of a familiar one. A second, stricter split retains only temperatures below 400~K during training (on the training materials) and evaluates exclusively on the 400--800~K range of the held-out test materials, isolating extrapolation performance under simultaneous material and temperature shift.

\section{Experimental Setup}

\subsection{Baselines}

We compare the proposed physics-guided model (PGML) against four baselines, all trained on the identical descriptor-plus-temperature input and the identical log-conductivity data loss:
\begin{itemize}
\item \textbf{Black-box MLP}: a five-layer feedforward network (hidden widths $128,128,64,32$) of comparable total parameter count to PGML's physics-parameter head, with no physics-derived structure and no shape-consistency penalty.
\item \textbf{Random Forest} and \textbf{Gradient-Boosted Trees}: standard scikit-learn tree-ensemble regressors (300 estimators; max depth 12 and 3, respectively), included as strong non-physics-aware reference points.
\item \textbf{Residual-only PRN} (physics-regularized network; we avoid the term PINN here since this baseline enforces a physical-consistency penalty rather than a PDE residual): the same Black-box MLP architecture, but trained with the shape-consistency penalty $\mathcal{L}_{\text{shape}}$ added to the loss (i.e., physics enforced only through the loss function, without the closed-form transport layer or physics-parameter head).
\end{itemize}

\subsection{Metrics}

We report: (i) RMSE and $R^2$ in $\log_{10}\sigma$ on the in-distribution test split; (ii) extrapolation RMSE on the 400--800~K held-out range; and (iii) a physical-consistency violation rate, defined as the fraction of test curves for which the model's predicted $d\sigma/dT$ has the wrong sign, relative to the regime expected from the material's known doping level, over more than 10\% of the evaluated temperature range. The consistency check applies equally to every model, tree ensembles included: Random Forest and Gradient-Boosted Trees are queried at every temperature on the same shared grid used for the neural models (not only the measurability-filtered points used for RMSE/$R^2$), and the same finite-difference Arrhenius-slope test is applied to their predicted curves -- a tree ensemble's \emph{training} objective has nothing to do with physics, but nothing prevents evaluating the physical plausibility of what it predicts at test time. Table~\ref{tab:main} reports this metric as not available for Random Forest and Gradient-Boosted Trees because that table's baseline configuration (fixed, untuned Gradient-Boosted Trees hyperparameters, matched to a since-superseded reduced closed form) is a historical result no longer reproducible from the released code (\texttt{code/README.md} explains why); every subsequent table reports it for both tree ensembles directly.

Paired significance testing across the 5 seeds (\texttt{code/significance\_tests.py}: paired $t$-test and Wilcoxon signed-rank) is reported throughout Section~5 wherever a numerical difference between models is discussed. With only 5 seeds, this testing has limited statistical power, and we treat it as exploratory rather than confirmatory evidence of an effect's absence or presence. As a complementary check less sensitive to any one seed's particular train/validation/test reshuffling, \texttt{code/bootstrap\_tests.py} additionally reports a material-level bootstrap for the paper's key comparisons: each model's per-curve RMSE is pooled across all 5 seeds' test sets, the benchmark's materials are resampled with replacement 10{,}000 times, and we report the resulting 95\% percentile confidence interval and two-sided bootstrap $p$-value for the difference in mean RMSE. This bootstrap resamples over the more natural unit of statistical independence here (materials, of which the benchmark has 24) rather than seeds (of which there are 5), but a given material can still appear as a test material in more than one seed, so it is not a bootstrap over fully independent units either; we accordingly report it as exploratory as well, alongside rather than in place of the paired seed-level tests.

\subsection{Training Details}

All neural models are trained full-batch (the entire training split fits comfortably in memory) with AdamW (learning rate $2\times10^{-3}$, weight decay $10^{-4}$, cosine decay over the schedule), for up to 2{,}000 epochs with early stopping on validation RMSE at a patience of 150 epochs. Tree ensembles in Tables~\ref{tab:main}--\ref{tab:richer} use 300 estimators (Random Forest: max depth 12; Gradient-Boosted Trees: max depth 3); Section~5.4 additionally tunes Gradient-Boosted Trees' hyperparameters via the validation split, and Table~\ref{tab:richer} onward uses that tuned version. Results are averaged over five random seeds controlling both weight initialization and the material-level train/validation/test split; we report mean $\pm$ standard deviation. With only 24 distinct materials, a material-level split of 70\%/15\%/15\% leaves 3--5 test materials per seed, so per-seed metrics -- particularly the consistency-violation rate, which is only defined on curves with an unambiguous intrinsic-dominated temperature window -- are individually noisy; we report the five-seed aggregate throughout and the total number of checkable curves pooled across seeds alongside it.

\section{Results}

All numbers below are produced by the accompanying code (\texttt{code/train.py}, \texttt{code/run\_multi\_seed.py}) and are exactly reproducible from the released synthetic benchmark; we report them as obtained, including where they run counter to our original hypothesis (Section~6 discusses why).

\subsection{Overall Accuracy and Extrapolation with a Reduced Closed Form}

We first report results with the reduced closed form described in an earlier iteration of Equation~\ref{eq:sigma_phys} -- a single combined exponential-times-power-law expression, $\sigma_{\text{phys}} \propto n_0 T^{3/2}\exp(-E_{g,\text{eff}}/2k_BT)\,\tau_0(T/T_0)^{-p}$, with four parameters per material rather than the five in the closed form we ultimately adopt -- because the comparison between the two closed forms is itself the paper's central finding (Section~5.2). Table~\ref{tab:main} summarizes five-seed performance across all models with this reduced closed form.

\begin{table}[h]
\centering
\caption{Test-set performance with the \emph{reduced} closed form, mean $\pm$ std over 5 seeds (each reshuffling both weight initialization and the material-level split). Lower is better for RMSE and violation rate; higher is better for $R^2$. Best result per column in bold. The consistency-violation rate is computed only on curves with an unambiguous intrinsic-dominated window; $n$ is the number of such curves pooled across all 5 seeds' test sets (out of $5\times$ 3--5 test materials $\times$ 3 doping levels each). Gradient-Boosted Trees here uses the fixed hyperparameters described in Section~4.3; Section~5.4 re-tunes them via validation, which changes this row's numbers in Tables~\ref{tab:richer}--\ref{tab:hybrid} slightly. Consistency violation is marked n/a for the tree ensembles in this table only: this table's specific baseline configuration (fixed, untuned Gradient-Boosted Trees hyperparameters, paired with the since-superseded reduced closed form) is a historical result no longer reproducible from the released code (\texttt{code/README.md}), so it was never computed for these two rows; Section~4.2 explains how the same check is computed for tree ensembles from Table~\ref{tab:richer} onward.}
\label{tab:main}
\begin{tabular}{lccccc}
\toprule
Model & RMSE ($\log_{10}\sigma$) & $R^2$ & Extrap.\ RMSE & Consistency violation (\%) & $n$ \\
\midrule
PGML, no RCN (physics-only) & $1.046 \pm 0.310$ & $0.744 \pm 0.138$ & $1.945 \pm 0.225$ & $18.3 \pm 15.3$ & 17 \\
PGML (full)                 & $1.048 \pm 0.300$ & $0.745 \pm 0.133$ & $1.944 \pm 0.249$ & $13.3 \pm 16.3$ & 17 \\
PGML, no $\mathcal{L}_{\text{shape}}$ (arch.\ only) & $0.892 \pm 0.350$ & $0.789 \pm 0.142$ & $1.368 \pm 0.440$ & $6.7 \pm 13.3$ & 17 \\
Black-box MLP               & $0.849 \pm 0.517$ & $0.748 \pm 0.357$ & $1.329 \pm 0.309$ & $\mathbf{0.0 \pm 0.0}$ & 17 \\
Random Forest                & $0.704 \pm 0.220$ & $0.891 \pm 0.062$ & $1.173 \pm 0.520$ & n/a\textsuperscript{*} & -- \\
\textbf{Residual-only PRN}  & $\mathbf{0.650 \pm 0.279}$ & $0.876 \pm 0.132$ & $1.175 \pm 0.582$ & $\mathbf{0.0 \pm 0.0}$ & 17 \\
Gradient-Boosted Trees      & $0.628 \pm 0.228$ & $\mathbf{0.908 \pm 0.057}$ & $\mathbf{1.307 \pm 0.549}$ & n/a\textsuperscript{*} & -- \\
\bottomrule
\end{tabular}
\end{table}

This result does not match our original hypothesis. The two PGML variants that use the reduced closed-form transport layer (full and physics-only) are the \emph{worst}-performing models on every column we can compute for them: highest RMSE, lowest $R^2$, worst extrapolation RMSE, and -- contrary to their entire motivation -- the highest physical-consistency violation rates in the table. The residual-only PRN (a plain Black-box MLP with the shape-consistency loss added, but no closed-form layer) is instead the strongest or joint-strongest model on RMSE, extrapolation, and consistency, matching the Black-box MLP's perfect 0\% violation rate while improving substantially on its accuracy. Gradient-Boosted Trees and Random Forest, which have no physics awareness at all, are also competitive with or better than both PGML variants on every metric.

\subsection{A Richer Closed Form Recovers Physical Consistency}

Rather than conclude from Table~\ref{tab:main} that architectural physics integration is unhelpful in general, we asked whether the reduced closed form itself was the problem: it commits to one effective activation energy and one effective scattering exponent per material, which cannot separately represent the extrinsic plateau and the intrinsic rise, nor two scattering channels with different temperature exponents (Section~3.1). We revised the closed-form transport layer to keep carrier concentration and mobility as separate terms with their own free parameters -- Equations~\ref{eq:n_closed}--\ref{eq:sigma_phys}, five parameters per material instead of four -- so that it sits in the same functional family as the benchmark's data-generating process rather than a further reduction of it, and re-ran the identical experiment (identical descriptors, splits, baselines, and evaluation protocol; only the PGML closed form changes, so the black-box and tree-ensemble rows are unchanged from Table~\ref{tab:main}; the residual-only PRN's numbers do shift slightly, since its shape-consistency loss shares the same $\lambda_1$ weight, which Section~3.4 explains was re-tuned for the richer closed form rather than held fixed). We report four PGML variants together with this richer base closed form: no high-temperature correction at all (\emph{no correction}), the RCN only, the explicit second conduction channel only (Equation~\ref{eq:second_channel}), and both together. Table~\ref{tab:richer} reports the outcome.

\begin{table}[h]
\centering
\caption{Test-set performance with the \emph{richer} base closed form (Equations~\ref{eq:n_closed}--\ref{eq:sigma_phys}) under four high-temperature-correction variants, mean $\pm$ std over 5 seeds. Baseline rows are identical to Table~\ref{tab:main} except Gradient-Boosted Trees, whose hyperparameters are tuned here on the validation split (Section~5.4) rather than fixed. Best result per column in bold.}
\label{tab:richer}
\begin{tabular}{lccccc}
\toprule
Model & RMSE ($\log_{10}\sigma$) & $R^2$ & Extrap.\ RMSE & Consistency violation (\%) & $n$ \\
\midrule
PGML (no correction)          & $0.861 \pm 0.304$ & $0.796 \pm 0.174$ & $1.253 \pm 0.465$ & $\mathbf{0.0 \pm 0.0}$ & 17 \\
PGML (RCN only)                & $0.889 \pm 0.292$ & $0.788 \pm 0.175$ & $1.549 \pm 0.503$ & $\mathbf{0.0 \pm 0.0}$ & 17 \\
PGML (2-channel)               & $0.797 \pm 0.288$ & $0.825 \pm 0.141$ & $1.304 \pm 0.487$ & $\mathbf{0.0 \pm 0.0}$ & 17 \\
PGML (2-channel + RCN)         & $0.769 \pm 0.224$ & $0.848 \pm 0.107$ & $1.637 \pm 0.465$ & $\mathbf{0.0 \pm 0.0}$ & 17 \\
Black-box MLP                  & $0.849 \pm 0.517$ & $0.748 \pm 0.357$ & $1.329 \pm 0.309$ & $\mathbf{0.0 \pm 0.0}$ & 17 \\
Random Forest                  & $0.704 \pm 0.220$ & $0.891 \pm 0.062$ & $1.173 \pm 0.520$ & $5.0 \pm 10.0$ & 17 \\
Residual-only PRN             & $0.765 \pm 0.306$ & $0.836 \pm 0.171$ & $\mathbf{1.118 \pm 0.306}$ & $\mathbf{0.0 \pm 0.0}$ & 17 \\
Gradient-Boosted Trees (tuned) & $\mathbf{0.656 \pm 0.196}$ & $\mathbf{0.909 \pm 0.042}$ & $1.166 \pm 0.422$ & $11.7 \pm 14.5$ & 17 \\
\bottomrule
\end{tabular}
\end{table}

The richer closed form changes the picture substantially. Every PGML variant now achieves a perfect $0\%$ consistency-violation rate, matching the Black-box MLP and residual-only PRN -- the physical-consistency problem in Table~\ref{tab:main} is resolved entirely, not merely reduced, once the closed form is expressive enough to represent the true carrier-concentration and mobility behavior without contorting a single exponent to fit both. The tree ensembles, evaluated on the same check for the first time in this table (Section~4.2), do \emph{not} share this perfect record: Random Forest violates on $5.0\% \pm 10.0\%$ of checkable curves and tuned Gradient-Boosted Trees on $11.7\% \pm 14.5\%$ -- noisy estimates at $n=17$ pooled curves, but a real, nonzero gap from every neural model in the table, despite the tree ensembles having no physics awareness and the best raw accuracy of anything tested. Section~6 returns to what this trade-off means for reading the tree ensembles' accuracy advantage as unqualified. Accuracy improves further with the explicit second channel: PGML (2-channel + RCN) reaches RMSE $0.769$ and $R^2$ $0.848$, matching the residual-only PRN's RMSE ($0.765$, within noise) and exceeding its $R^2$ ($0.836$) -- the best result any PGML variant achieves in this study, and the first to be competitive with, rather than clearly behind, the strongest neural baseline. It does not, however, close the gap to the tree ensembles ($0.656$--$0.704$ RMSE with Gradient-Boosted Trees' hyperparameters tuned on validation, Section~5.4), and the plain second-channel variant without the RCN (\emph{2-channel}, RMSE $0.797$) generalizes to unseen materials at unseen temperatures better than the combined variant (extrapolation RMSE $1.304$ vs.\ $1.637$) -- adding the RCN on top of an already-good closed form appears to buy a little in-distribution accuracy at the cost of extrapolation robustness. Section~6 discusses this trade-off together with the ablation below.

\subsection{Ablation: The Second Conduction Channel versus the Residual Correction Network}

Table~\ref{tab:ablation} isolates the effect of the two mechanisms for handling the narrow-gap high-temperature correction the base closed form omits (same data as Table~\ref{tab:richer}, reorganized for the ablation comparison). Consistency violation is $0.0\% \pm 0.0\%$ for all four variants and is omitted as a column.

\begin{table}[h]
\centering
\caption{Ablation of the second conduction channel and the RCN on top of the richer base closed form, mean $\pm$ std over 5 seeds.}
\label{tab:ablation}
\begin{tabular}{lcc}
\toprule
Variant & RMSE ($\log_{10}\sigma$) & Extrap.\ RMSE \\
\midrule
PGML (no correction)     & $0.861 \pm 0.304$ & $\mathbf{1.253 \pm 0.465}$ \\
PGML (RCN only)          & $0.889 \pm 0.292$ & $1.549 \pm 0.503$ \\
PGML (2-channel)         & $0.797 \pm 0.288$ & $1.304 \pm 0.487$ \\
PGML (2-channel + RCN)   & $\mathbf{0.769 \pm 0.224}$ & $1.637 \pm 0.465$ \\
\bottomrule
\end{tabular}
\end{table}

Physical consistency is no longer the axis that distinguishes these variants -- all four achieve $0.0\% \pm 0.0\%$ violation -- so the ablation speaks only to accuracy and extrapolation, and the two-by-two structure (channel $\times$ RCN, on or off) separates their contributions. The explicit second channel is a clear net positive for accuracy on its own: \emph{2-channel} improves RMSE over \emph{no correction} ($0.861 \to 0.797$) while also improving extrapolation ($1.253 \to 1.304$ is within noise of each other, both far better than \emph{RCN only}'s $1.549$). The RCN's effect depends on what it is added to: added to \emph{no correction} it barely moves RMSE and makes extrapolation markedly worse (\emph{RCN only}, $1.549$); added on top of the second channel it gives the best RMSE of all four variants (\emph{2-channel + RCN}, $0.769$) but the worst extrapolation ($1.637$). We read this as the RCN doing two different jobs depending on the closed form under it: when the closed form is weak (no correction), the RCN has real structure to learn and generalizes poorly because it is learning the wrong (unstructured) thing; when the closed form is already strong (2-channel), the RCN has only noise-level residual error left to fit, and fitting that noise on 51 training curves is what degrades extrapolation to materials that were never part of that noise. Section~6 discusses the resulting accuracy/extrapolation trade-off between \emph{2-channel} and \emph{2-channel + RCN} in more detail.

\subsection{Attempting to Close the Gap to Tree Ensembles: Hybrids and Hyperparameter Tuning}

Tables~\ref{tab:main}--\ref{tab:richer} leave Random Forest and Gradient-Boosted Trees as the most accurate models throughout, un-moved by either closed-form revision. We tried three things to test whether that gap was closable at all with the tools in this paper. First, since the Gradient-Boosted Trees baseline throughout Tables~\ref{tab:main}--\ref{tab:richer} used a single fixed hyperparameter setting (300 estimators, max depth 3) rather than a tuned one, we ran a small grid search (36 combinations of estimator count, depth, learning rate, and minimum leaf size) selected by RMSE on the material-level validation split, separately for the standard and the T$\leq$400K extrapolation-training regime, and refit on the training split alone with the selected hyperparameters -- isolating the effect of tuning from any change in training data. Second and third, we tried two hybrids between the tuned Gradient-Boosted Trees and the trained \emph{PGML (2-channel)} model (no RCN, chosen as the best-extrapolating richer-closed-form variant, Table~\ref{tab:richer}). The first hybrid keeps the physics prediction and lets a tree ensemble fit whatever it does not capture: we compute the closed form's prediction on every training point, fit a Gradient-Boosted Trees regressor (same tuned hyperparameters) on the \emph{residual} between the observed and the physics-predicted $\log_{10}\sigma$ using the same descriptor-plus-temperature features as the tree baselines, and add the two predictions at test time. The second instead hands the tree ensemble the closed form's fitted physics parameters ($E_{g,\text{eff}}, C_n, N_0, \tau_{\text{ac}0}, \tau_{\text{ii}0}$, and the second-channel parameters) as additional engineered features alongside the raw descriptors and $T$, and lets it predict $\log_{10}\sigma$ directly rather than a residual. Table~\ref{tab:hybrid} reports all three outcomes, alongside the two closed-form-only models and Random Forest (left untuned) for reference.

\begin{table}[h]
\centering
\caption{Hyperparameter-tuned Gradient-Boosted Trees and two physics-plus-boosting hybrids built on it, compared to their components, mean $\pm$ std over 5 seeds. Best result per column in bold.}
\label{tab:hybrid}
\begin{tabular}{lccc}
\toprule
Model & RMSE ($\log_{10}\sigma$) & $R^2$ & Extrap.\ RMSE \\
\midrule
PGML (2-channel)                    & $0.797 \pm 0.288$ & $0.825 \pm 0.141$ & $1.304 \pm 0.487$ \\
Gradient-Boosted Trees (tuned)       & $\mathbf{0.656 \pm 0.196}$ & $\mathbf{0.909 \pm 0.042}$ & $1.166 \pm 0.422$ \\
Random Forest (untuned)              & $0.704 \pm 0.220$ & $0.891 \pm 0.062$ & $1.173 \pm 0.520$ \\
PGML (2-channel) + GBT residual      & $0.747 \pm 0.272$ & $0.846 \pm 0.137$ & $\mathbf{1.100 \pm 0.334}$ \\
GBT + physics-parameter features     & $0.775 \pm 0.267$ & $0.855 \pm 0.095$ & $1.415 \pm 0.399$ \\
\bottomrule
\end{tabular}
\end{table}

None of the three attempts closes the gap, and each tells a different part of the story. Tuning Gradient-Boosted Trees' own hyperparameters barely moves its numbers (RMSE $0.628 \to 0.656$, $R^2$ $0.908 \to 0.909$, both changes far smaller than the five-seed standard deviation; extrapolation RMSE improves somewhat, $1.307 \to 1.166$) -- the fixed hyperparameters used throughout Tables~\ref{tab:main}--\ref{tab:richer} were already close to as good as a 36-point validation-selected grid search finds on this benchmark, so the tree ensembles' advantage is not an artifact of under-tuning the baseline. The two hybrids, now built on this tuned Gradient-Boosted Trees, diverge instructively from each other. The residual hybrid's RMSE ($0.747$) and $R^2$ ($0.846$) sit between the physics-only model and Gradient-Boosted Trees, an improvement over the physics prior alone but short of matching the tree ensemble it is partly built from; it also achieves the best extrapolation RMSE of every model in this paper on this synthetic benchmark, physics-guided or not ($1.100$, versus $1.166$--$1.637$ for every other row across Tables~\ref{tab:main}, \ref{tab:richer}, and this one) -- better than the physics-only model it starts from, and better than the tuned Gradient-Boosted Trees model supplying its residual correction. The feature-augmentation hybrid, by contrast, is not an improvement on anything: its RMSE ($0.775$) and $R^2$ ($0.855$) are both worse than tuned Gradient-Boosted Trees, and its extrapolation RMSE ($1.415$) is worse than both the residual hybrid's and plain Gradient-Boosted Trees'. Handing the tree ensemble the closed form's fitted parameters bought it nothing, and mildly hurt it, likely because those parameters are a deterministic (if nonlinear) function of the same descriptors the tree ensemble already receives -- extra, correlated dimensions for a boosted tree to split on, without new information, on a dataset small enough (51 training curves) that the added dimensionality is itself a cost. We read this three-way contrast as evidence that the tree ensembles' advantage on this benchmark is a genuine, tuning-robust property of the learning problem rather than an artifact of an unfairly weak baseline, and that what makes the residual hybrid work is specifically the closed form's \emph{output} value (a good prior estimate of $\log_{10}\sigma$ itself) rather than its internal parameters or the flexibility of the tree ensemble correcting it: the physics prior and the tree-based residual are contributing complementary things only when the tree ensemble is given the prior's prediction to correct, not merely the ingredients that produced it or a better-tuned version of itself. None of these three interventions treats the problem as one of \emph{integration} across every model tried so far, rather than a single best replacement for the closed form; Section~5.5 tries that directly.

\subsection{An Integrated Ensemble System}

Rather than one more single model, we tried combining every model in this study -- the four PGML variants, the Black-box MLP, the residual-only PRN, Random Forest, and tuned Gradient-Boosted Trees, eight base predictors in all -- into one final predictor, two ways. The first is a simple unweighted average of all eight base predictions. The second is a stacked ensemble: a non-negative-weight Ridge meta-learner fits combination weights on the validation split's base-model predictions (data none of the eight base models' own parameters were fit on) and applies those weights to the test-set base predictions, separately for the standard and extrapolation regimes. Table~\ref{tab:ensemble} reports both, alongside the strongest individual models from Tables~\ref{tab:richer}--\ref{tab:hybrid} for reference.

\begin{table}[h]
\centering
\caption{Integrating all eight base models into one predictor, mean $\pm$ std over 5 seeds. Best result per column in bold.}
\label{tab:ensemble}
\begin{tabular}{lccc}
\toprule
Model & RMSE ($\log_{10}\sigma$) & $R^2$ & Extrap.\ RMSE \\
\midrule
Gradient-Boosted Trees (tuned)       & $0.656 \pm 0.196$ & $\mathbf{0.909 \pm 0.042}$ & $1.166 \pm 0.422$ \\
Random Forest (untuned)              & $0.704 \pm 0.220$ & $0.891 \pm 0.062$ & $1.173 \pm 0.520$ \\
PGML (2-channel) + GBT residual      & $0.747 \pm 0.272$ & $0.846 \pm 0.137$ & $1.100 \pm 0.334$ \\
\textbf{Ensemble (simple average)}   & $\mathbf{0.630 \pm 0.197}$ & $0.896 \pm 0.084$ & $\mathbf{0.937 \pm 0.297}$ \\
Ensemble (stacked, Ridge)            & $0.648 \pm 0.171$ & $0.895 \pm 0.073$ & $1.121 \pm 0.347$ \\
\bottomrule
\end{tabular}
\end{table}

The simple average has, by a modest margin, the lowest mean RMSE of any model or combination in the entire study ($0.630$, versus $0.656$ for tuned Gradient-Boosted Trees), and by a clearer margin the lowest mean extrapolation RMSE ($0.937$, versus $1.100$ for the next-best row, the residual hybrid, and $1.166$--$1.637$ for everything else across Tables~\ref{tab:main}--\ref{tab:hybrid}). Its $R^2$ ($0.896$) is close to but still short of tuned Gradient-Boosted Trees' ($0.909$). None of these differences from Gradient-Boosted Trees reaches significance in a paired test across the 5 seeds, nor in the exploratory material-level bootstrap of Section~4.2 (pooled RMSE $0.657$ vs.\ $0.683$, 95\% CI $[-0.300, +0.269]$, $p = 0.82$; Section~5.7 reports both exploratory tests' numbers in full on the stronger, weighted variant of this ensemble, where the unweighted average's differences are of comparable or smaller magnitude), so we read this as the closest this paper comes to \emph{numerically} narrowing the gap -- not a better closed form, not a better tree ensemble, but combining every model tried, none of which individually closed it, into one predictor whose mean sits at or ahead of the strongest individual baseline on two of three metrics, though not distinguishably so given only 5 seeds, and within noise on the third. Because the combination includes Gradient-Boosted Trees itself, this is also not independent evidence that architectural physics integration outperforms it -- Section~6 discusses this distinction.

The stacked ensemble does not improve on the simple average, despite having strictly more information available to it (validation-selected weights rather than a fixed uniform one) -- its RMSE ($0.648$) and extrapolation RMSE ($1.121$) are both worse. Inspecting the per-seed learned weights explains why: they are highly unstable. Across the five seeds, the Ridge meta-learner variously puts nearly all its weight on Gradient-Boosted Trees and Random Forest (one seed), on the four PGML variants roughly equally (another), or on the Black-box MLP alone (a third, with a weight of $0.95$) -- different seeds' 12 validation curves ($\sim$300 points) apparently do not contain enough signal to reliably distinguish which of eight correlated base models will generalize best to the test materials, so the meta-learner overfits to validation-specific noise in a different direction each time. The unweighted average sidesteps this entirely by not trying to estimate weights from data at all, which on a benchmark this small is itself a form of regularization. A natural next question is whether the ensemble itself can be improved by widening its pool of candidates -- Section~5.6 tries that directly.

\subsection{Extending the Ensemble with the Physics-Plus-Boosting Hybrids}

Both ensembles in Table~\ref{tab:ensemble} draw only on the eight models trained for Tables~\ref{tab:richer} and \ref{tab:hybrid}; neither the residual hybrid nor the feature-augmentation hybrid from Section~5.4 is a candidate member, even though both are already trained by that point and cost nothing extra to add. A sixth experiment repeats both combination strategies -- unweighted average and stacked Ridge -- over all ten models (the original eight plus both hybrids), to test whether a bigger candidate pool closes more of the remaining gap. Table~\ref{tab:ensemble2} reports the result alongside the original two rows for comparison.

\begin{table}[h]
\centering
\caption{Extending both ensemble strategies with the two physics-plus-boosting hybrids as additional candidate members (ten models total), mean $\pm$ std over 5 seeds. Best result per column in bold.}
\label{tab:ensemble2}
\begin{tabular}{lccc}
\toprule
Model & RMSE ($\log_{10}\sigma$) & $R^2$ & Extrap.\ RMSE \\
\midrule
\textbf{Ensemble (8 models, simple average)}   & $\mathbf{0.630 \pm 0.197}$ & $0.896 \pm 0.084$ & $\mathbf{0.937 \pm 0.297}$ \\
Ensemble (8 models, stacked Ridge)             & $0.648 \pm 0.171$ & $0.895 \pm 0.073$ & $1.121 \pm 0.347$ \\
Ensemble (10 models, simple average)           & $0.634 \pm 0.189$ & $\mathbf{0.896 \pm 0.079}$ & $0.964 \pm 0.304$ \\
Ensemble (10 models, stacked Ridge)            & $0.645 \pm 0.179$ & $0.894 \pm 0.075$ & $1.036 \pm 0.260$ \\
\bottomrule
\end{tabular}
\end{table}

Adding both hybrids as members changes neither combination strategy in a way that matters. The ten-model simple average ($0.634$ RMSE, $0.896$ $R^2$, $0.964$ extrapolation RMSE) is not distinguishable from the original eight-model average on either metric in a paired test across the 5 seeds ($p = 0.75$ RMSE, $p = 0.28$ extrapolation RMSE), so this reads as noise rather than a real degradation or improvement. This is a mildly informative null result on its own: since both hybrids are individually weaker than tuned Gradient-Boosted Trees (Table~\ref{tab:hybrid}), one might expect diluting a strong average with two more mediocre predictors to hurt it more visibly; that it does not suggests the simple average's variance-reduction mechanism is not sensitive to exactly which reasonable models are included, at least at this pool size. The ten-model stacked ensemble ($0.645$ RMSE, $1.036$ extrapolation RMSE) is likewise not distinguishable from its eight-model counterpart ($p = 0.77$ RMSE, $p = 0.38$ extrapolation RMSE), and the same instability mechanism identified in Section~5.5 persists: the Ridge meta-learner still has only 12 validation curves to estimate weights over, now for ten correlated candidates instead of eight, so giving it more choices does not give it more information to choose well with. Neither combination strategy's ranking changes -- the plain average remains ahead of the stack regardless of how many models it is allowed to draw on. We read this as evidence that the ensemble's earlier result (Section~5.5) was not leaving easy gains on the table by excluding the hybrids specifically; the limiting factor is the combination strategy and the size of the validation split available to fit one, not which models are eligible to be combined -- a claim Section~5.7 tests directly by trying a third combination strategy, between the two extremes of no data-driven weighting and a fully free regression fit.

\subsection{Weighting the Ensemble by Validation Accuracy}

The simple average (Section~5.5) uses no validation-set information to combine models; the stacked Ridge uses a full regression fit, with as many free weights as base models, on the same 12 validation curves, and that many free parameters is exactly what made it unstable. A seventh experiment tries a middle ground: weight each base model's prediction by the inverse of its own \emph{squared} validation RMSE (inverse-variance weighting), normalized to sum to one. This uses the validation split to estimate exactly one scalar per model -- how accurate it is on held-out data -- rather than a full covariance-aware regression over correlated candidates, and is the classical minimum-variance combination rule for averaging independent unbiased estimators. We compute it for both the eight-model and ten-model pools from Sections~5.5--5.6. Table~\ref{tab:ensemble3} reports the result.

\begin{table}[h]
\centering
\caption{Weighting the ensemble by inverse-variance validation RMSE instead of an unweighted average or a freely-fit stack, mean $\pm$ std over 5 seeds. Best result per column in bold.}
\label{tab:ensemble3}
\resizebox{\textwidth}{!}{%
\begin{tabular}{lccc}
\toprule
Model & RMSE ($\log_{10}\sigma$) & $R^2$ & Extrap.\ RMSE \\
\midrule
Gradient-Boosted Trees (tuned)                     & $0.656 \pm 0.196$ & $\mathbf{0.909 \pm 0.042}$ & $1.166 \pm 0.422$ \\
Ensemble (8, simple average)                       & $0.630 \pm 0.197$ & $0.896 \pm 0.084$ & $0.937 \pm 0.297$ \\
Ensemble (8, stacked Ridge)                        & $0.648 \pm 0.171$ & $0.895 \pm 0.073$ & $1.121 \pm 0.347$ \\
\textbf{Ensemble (8, weighted by val.\ RMSE)}      & $\mathbf{0.593 \pm 0.202}$ & $0.905 \pm 0.083$ & $\mathbf{0.822 \pm 0.258}$ \\
Ensemble (10, weighted by val.\ RMSE)              & $0.604 \pm 0.200$ & $0.902 \pm 0.082$ & $0.834 \pm 0.289$ \\
\bottomrule
\end{tabular}%
}
\end{table}

Validation-RMSE weighting has a numerically lower mean on every metric than both earlier combination strategies. Against the simple average: RMSE $0.630 \to 0.593$, $R^2$ $0.896 \to 0.905$, extrapolation RMSE $0.937 \to 0.822$; a paired $t$-test across the 5 seeds supports the RMSE improvement ($p = 0.012$; Wilcoxon signed-rank $p = 0.063$, the smallest attainable at $n=5$), but not the extrapolation-RMSE improvement ($p = 0.12$). The exploratory material-level bootstrap (Section~4.2) agrees on the RMSE improvement: pooling every test curve's squared error across all 5 seeds and resampling the 12 materials that appear in at least one seed's test set, weighted-RMSE averaging's pooled RMSE ($0.624$) is lower than the simple average's ($0.657$), a difference whose 95\% CI ($[-0.057, -0.008]$) excludes zero ($p = 0.008$) -- convergent evidence, from a resampling axis independent of the paired seed-level test, that this specific improvement is real rather than an artifact of which 5 splits happened to be drawn. Against the stacked Ridge, every metric is numerically better but neither the paired test ($p = 0.48$ RMSE, $p = 0.15$ extrapolation RMSE) nor the bootstrap (pooled RMSE $0.624$ vs.\ $0.666$, 95\% CI $[-0.151, +0.053]$, $p = 0.40$) supports significance. Against tuned Gradient-Boosted Trees, the previous strongest single baseline: mean RMSE is lower ($0.593$ vs.\ $0.656$) and mean extrapolation RMSE is lower ($0.822$ vs.\ $1.166$), with $R^2$ within noise of it ($0.905 \pm 0.083$ vs.\ $0.909 \pm 0.042$) -- but neither a paired $t$-test across the same 5 seeds ($p = 0.70$ RMSE, $p = 0.93$ $R^2$, $p = 0.19$ extrapolation RMSE) nor the material-level bootstrap (pooled RMSE $0.624$ vs.\ $0.683$, 95\% CI $[-0.325, +0.225]$, $p = 0.66$) supports any of these three differences as statistically significant; with only 5 seeds or 12 bootstrappable materials, both comparisons have limited power, and the numerical gap should be read as suggestive rather than established. It is, in any case, not independent evidence that the physics-guided architecture outperforms Gradient-Boosted Trees, since the ensemble being compared \emph{includes} Gradient-Boosted Trees as one of its members -- Section~6 returns to this point. (The bootstrap's pooled RMSE values differ slightly from the paired tests' seed-averaged means above -- e.g. $0.624$ vs.\ $0.593$ for the weighted ensemble -- because pooling every curve's squared error across all 5 seeds into one RMSE is a different statistic from averaging 5 separately-computed per-seed RMSEs; the two are complementary, not literally the same quantity computed twice.) The ten-model version (adding both hybrids, as in Section~5.6) performs almost identically ($0.604$ RMSE, $0.834$ extrapolation RMSE) -- consistent with Section~5.6's finding that pool size is not the limiting factor for a combination strategy that is already working well, only for one that is not.

Comparing the per-seed weights this scheme assigns (recorded as \texttt{validation\_weights} in \texttt{code/results\_seed*.json}) to the stacked Ridge's weights explains the improvement. The stacked Ridge, across five seeds, ranged from putting $0.95$ of its weight on the Black-box MLP alone to putting essentially all of it on Gradient-Boosted Trees and Random Forest -- a mean per-model weight standard deviation, across seeds, of $0.140$. Validation-RMSE weighting is far steadier: every model gets a nontrivial share in every seed (no weight ever drops below roughly $0.03$ or exceeds roughly $0.22$ in the eight-model pool), for a mean per-model weight standard deviation of $0.047$, roughly a third as much. This is the direct consequence of fitting one number per model instead of a joint regression over eight to ten correlated predictors on the same 12 validation curves: inverse-variance weighting cannot overfit the interactions between candidates because it never looks at them, only at each model's own marginal accuracy. That turns out to matter more than the extra flexibility a free regression fit nominally provides -- the closest this paper comes to a complete answer on how to combine models when no single one closes a performance gap on a small benchmark.

\section{Discussion}

The comparison between Table~\ref{tab:main} and Table~\ref{tab:richer} is the central finding of this study, and it is a comparison we would not have made had the first (reduced-form) experiment happened to look reasonable. Routing predictions through a closed-form transport layer did not, by itself, improve accuracy, extrapolation, or physical consistency relative to an unconstrained network on this benchmark -- with the reduced closed form it made all three markedly worse, while the identical shape-consistency loss term helped when applied to an unconstrained network. Revising only the closed form's functional richness, with every other part of the architecture, loss, training procedure, and benchmark held fixed, resolved the physical-consistency failure completely and closed most of the accuracy and extrapolation gap. This isolates the reduced closed form's expressive power, not the general idea of architectural physics integration, as the cause of the original negative result: Equation~\ref{eq:sigma_phys}'s earlier four-parameter form committed to a single combined exponential-times-power-law shape per material, which cannot separately represent the extrinsic plateau and the intrinsic rise, nor two scattering channels with different temperature exponents, while the benchmark's generator (Section~3.5) blends these through a smooth square-root transition and Matthiessen's rule with independent exponents. Across the crossover region between the extrinsic and intrinsic regimes -- exactly the region a temperature-dependent conductivity model is most useful for -- the reduced closed form was a biased approximation, and that bias was not something more training data, a better optimizer, or a larger residual correction could remove: a $\pm 15\%$ bounded RCN barely moved any metric in Table~\ref{tab:main}, and loosening the bound to $\pm 50\%$ in preliminary experiments (not shown) did not help either, because the error was in the closed form's functional shape rather than its amplitude. The richer closed form (Equations~\ref{eq:n_closed}--\ref{eq:sigma_phys}) removes that bias by construction: it is in the same functional family as the generator, so the physics-parameter head only has to identify the right five numbers per material rather than compensate for a missing degree of freedom.

That the richer base closed form still trailed the residual-only PRN on RMSE (Table~\ref{tab:richer}, \emph{no correction} and \emph{RCN only} rows) motivated one further revision: giving the narrow-gap high-temperature effect (Section~3.5) an explicit, structurally matched term rather than leaving it to a generic residual network. This closed the remaining gap almost entirely. \emph{PGML (2-channel + RCN)} reaches RMSE $0.769$, matching the residual-only PRN's $0.765$ within noise, and its $R^2$ of $0.848$ exceeds the residual-only PRN's $0.836$ -- the best result any PGML variant achieves in this study, and, together with Table~\ref{tab:richer}'s consistency-violation column, the first PGML variant to be competitive with the strongest neural baseline on every accuracy metric simultaneously rather than trading one for another. The mechanism is the same one that made the base closed-form revision work: the explicit second channel commits to the generator's actual functional form (a single logistic turn-on, Equation~\ref{eq:second_channel}) rather than an unconstrained correction of arbitrary shape, so the physics-parameter head only has to identify three more material-specific numbers ($\gamma_2$, $T_{c2}$, $T_{w2}$) rather than have a separate network rediscover that the correction should look like a sigmoid in the first place.

The remaining gap to the tree ensembles (Table~\ref{tab:richer}, $0.656$--$0.704$ RMSE) did not close, and we do not have a structural explanation of the same kind for it: Random Forest and Gradient-Boosted Trees have no notion of the physics at all, and their advantage is presumably the generic one tree ensembles often have on small, tabular, non-smooth regression problems -- documented well outside materials science, e.g.\ in Grinsztajn et al.'s benchmark of 45 tabular datasets, where tree-based models remain state of the art on medium-sized tabular data regardless of neural architecture \citep{grinsztajn2022tree} -- not something a better closed form would be expected to fix. We flag this explicitly because it would be easy to over-read the second channel's success as ``physics-guided modeling wins once you get the closed form right'' -- on this benchmark, with 51 training curves and a 9-dimensional descriptor vector, Gradient-Boosted Trees were the single best model throughout every experiment in this paper, physics-guided or not. That advantage, however, is not free of a physical-consistency cost: once the same check is applied to the tree ensembles' predictions (Table~\ref{tab:richer}, Section~4.2), Random Forest and Gradient-Boosted Trees violate the expected Arrhenius-slope sign on $5.0\%$ and $11.7\%$ of checkable curves respectively, versus a perfect $0\%$ for every neural model in the table, physics-guided or not. Nothing in a tree ensemble's training objective discourages a predicted curve from locally decreasing with temperature in a region where the physics says it should rise, or vice versa, and apparently this happens often enough, on a benchmark this size, to show up even in a metric the tree ensembles were never optimized against. The practical upshot is that ``Gradient-Boosted Trees is simply the best model here'' is only true along the accuracy and extrapolation axes; along the physical-plausibility axis every architecture with an explicit physics-consistency mechanism -- PGML, the residual-only PRN, and, incidentally, the unconstrained Black-box MLP too -- outperforms it, which is exactly the kind of trade-off a raw accuracy leaderboard would hide.

Section~5.4's three interventions were a direct test of that claim from three angles: if the tree ensembles' advantage were an artifact of an under-tuned baseline, a validation-selected hyperparameter grid should have closed most of it; if it were a generic tabular-regression effect entirely separable from the smooth part of the signal, handing a tree ensemble the residual on top of the best closed form should have closed most of the rest. Neither happened. Tuning barely moved Gradient-Boosted Trees' RMSE ($0.628 \to 0.656$, within one-third of a standard deviation), confirming the gap in Tables~\ref{tab:main}--\ref{tab:richer} was not simply a matter of unfair defaults. The residual hybrid closed some of what remained ($0.797 \to 0.747$ RMSE, roughly a third of the way from the physics-only model to tuned Gradient-Boosted Trees) but not all, which we read as mixed evidence: the hybrid's residual-fitting problem is presumably a little easier than the tree ensemble's original problem (the closed form has already removed the smooth, physically-structured part of the signal), yet the boosted correction still falls short of what Gradient-Boosted Trees achieves when given the whole signal directly, suggesting the tree ensemble's advantage is not purely about the smooth part of the function that the closed form already captures. The clearest result from the hybrid is not on the RMSE axis at all: it achieved the best extrapolation RMSE of every model in the paper on this synthetic benchmark ($1.100$), better than both of its own components. Our reading is that the closed form's main practical value is not raw in-distribution accuracy -- where a strong tabular learner given enough data, and no more hyperparameter tuning than the closed form itself received, can apparently do as well or better without any physics -- but the behavior of the resulting function outside the training distribution, which a tree ensemble has no mechanism to get right beyond what its training leaves fall inside of, and which a boosted correction \emph{on top of} a sensible closed form does not seem to damage. This suggests hybrids of this kind may be more promising as an extrapolation-robustness technique than as a pure accuracy technique, which is a different framing than the one usually offered for architectural physics integration.

Section~5.5's simple-average ensemble goes one step further, numerically if not statistically: $0.630$ mean RMSE edges out tuned Gradient-Boosted Trees' $0.656$, and $0.937$ mean extrapolation RMSE clears every other row in the paper by a wider margin than any single revision managed, though (Section~5.7 reports the paired test) neither difference is distinguishable from noise at $n=5$ seeds. We do not think this contradicts the reading above -- it is consistent with it. An unweighted average is, in effect, a crude variance-reduction device: each of the eight base models makes different errors on different test materials (the closed-form models err where their functional assumptions are weakest, the tree ensembles err where the training set has too little local support, the black-box network and the PRN somewhere in between), and averaging cancels some of that idiosyncratic error without requiring any model to be individually correct everywhere. That this beats the single best model at all is itself informative about how much of the remaining gap in Tables~\ref{tab:richer}--\ref{tab:hybrid} was individual-model variance rather than a systematic deficiency of every architecture that is not a tree ensemble. It is also why the stacked variant, which tries to do better than uniform weighting by estimating which models to trust, underperforms the naive average: with only 12 validation curves, the weight-estimation problem itself has high variance, and the learned weights (Section~5.5) chase validation-set noise in a different direction each seed. The lesson we take is not that ensembling is a free way to beat tree ensembles in general, but that on a benchmark this small, a method that avoids estimating anything extra from data (uniform averaging) can outperform one that tries to be smarter about it (learned weighting) -- the opposite of what more data would likely show, and a caution against reading the stacked ensemble's underperformance as evidence that learned combination is a bad idea rather than a data-starved one here.

Sections~5.6 and 5.7 push this further and sharpen what ``closing the gap'' means here. Extending the ensemble to include both physics-plus-boosting hybrids left the simple average and the stacked Ridge essentially unchanged, so the earlier ensemble's success was not an artifact of one particular eight-model roster; the bottleneck was the combination rule and the small validation split, not which candidates were on offer. Weighting each model's prediction by the inverse of its own squared validation RMSE, rather than averaging uniformly or fitting a fully free regression, then closed the remaining gap on every metric at once, including $R^2$, where the earlier ensembles still trailed tuned Gradient-Boosted Trees. We want to be precise about what this does and does not show. The validation-RMSE-weighted ensemble's advantage comes from combining every model tried, the two tree ensembles included, not from the physics-guided architecture outperforming them unaided -- no PGML variant, in Sections~5.1--5.3, ever matched the tree ensembles' accuracy on its own, and nothing in Sections~5.6--5.7 changes that. What these two experiments show instead is that once a reasonably diverse set of models is available, a physics-guided one among them, the choice of how to combine them can matter as much as which individual model is strongest, and that a combination rule cheap enough not to overfit a 12-curve validation split can outperform one that tries to be cleverer about it. Whether a physics-guided model can win outright on this benchmark, rather than by being folded into a mixture with the very baselines it is being compared against, remains an open question these two experiments do not answer.

The ablation (Table~\ref{tab:ablation}) also surfaces an accuracy/extrapolation trade-off that repeats a pattern from Section~5.1: \emph{2-channel + RCN} has the best in-distribution accuracy but the worst extrapolation RMSE among the four richer-closed-form variants ($1.637$, worse even than \emph{RCN only}'s $1.549$), while the plain \emph{2-channel} variant without the RCN has both good accuracy ($0.797$) and the best extrapolation apart from \emph{no correction} ($1.304$). We think this is the RCN doing different jobs in different contexts. Added on top of a weak closed form (\emph{no correction}), the RCN has genuine structure to learn (the entire narrow-gap correction) and generalizes it poorly to unseen materials, because an unconstrained network's approximation of a sigmoid turn-on does not extrapolate the way an actual sigmoid does. Added on top of an already-strong closed form (\emph{2-channel}), the RCN has little genuine signal left and mostly fits residual noise specific to the 51 training curves, which is precisely the kind of fitting that does not transfer to new materials at new temperatures. In both cases the RCN's lack of a committed functional form -- its main advantage in principle, since it can absorb any residual effect -- is also what limits how well what it learns generalizes.

We want to be careful about how far this generalizes. It is a demonstration that, on \emph{this particular} synthetic generator, whether architectural physics integration helps, hurts, or does neither depends on a design choice -- the closed form's functional richness -- that is easy to get wrong and easy not to test, and that testing it against both a black-box baseline and a loss-only physics-informed baseline is what exposed the problem, and iterating on the closed form (first its base transport terms, then its high-temperature correction) is what closed most of the resulting gap. It is not evidence that an eight-parameter closed form is the correct or final answer, nor that architectural physics integration will always eventually catch up to a flexible baseline if you iterate on the closed form enough -- the tree ensembles' advantage on this benchmark did not close, and we have no principled account of why a closed form should be expected to close it. The practical lesson we draw is methodological, and it cuts in both directions: a claimed advantage of architectural physics integration should be checked against both a black-box and a loss-only physics-informed baseline before being taken as general (it was absent for our first closed form), but a claimed \emph{disadvantage} should equally be checked against a more expressive closed form before concluding that architectural integration itself, rather than one specific choice of closed form, is the problem (it was, in our case, a fixable choice rather than a fundamental one).

\section{Limitations}
\label{sec:limitations}

Five main limitations bound the interpretation of these results; secondary details behind each are grouped below rather than given as separate points.

\textbf{1. Synthetic data and a small, literature-parameterized benchmark.} All $\sigma(T)$ data are synthetic, generated by a reduced transport model (Section~3.5) rather than measured, and most of its 9 electronic-structure descriptors (Section~3.2) are representative literature-typical values rather than original per-material calculations. The richer base closed form and the explicit second channel are deliberately in the same functional family as this generator, which is precisely what let them recover physical consistency and close most of the accuracy gap -- but this also means our results are, in part, a demonstration that a model fits data well when its functional form matches the process that produced it, not general evidence about real, measured $\sigma(T)$ curves. Separately, the benchmark spans only 24 materials, leaving 3--5 held out per seed at a material-level 70/15/15 split; the paired significance tests we do report (Sections~5.5--5.7) have correspondingly limited statistical power at $n=5$ seeds, and standard deviations are often large relative to the gaps between closer-performing models, so specific orderings should be read as suggestive rather than tightly estimated even where the qualitative pattern (richer beats reduced, second channel beats RCN alone) is consistent across all five seeds. The material-level bootstrap we report alongside these tests (Section~4.2) resamples a larger and arguably more natural unit (materials rather than seeds), but the same 24-material ceiling limits it too -- at most 12 of the 24 materials appear as a test material in any of the 5 seeds' splits, and a material appearing in more than one seed's test set is not resampled as an independent unit either -- so we report both kinds of test as exploratory rather than confirmatory throughout.

\textbf{2. Hyperparameters and loss weights chosen once, not tuned per variant.} The shape-consistency loss's finite-difference slope estimate (25-point grid) is somewhat noisy and was not re-ablated against the final architecture; its weight $\lambda_1$, the RCN bound $\eta$, and the second channel's parameter ranges were each set from a small manual sweep on one seed rather than tuned independently per variant. Section~5.1's finding that $\lambda_1$'s optimum depends on the base closed form's own accuracy should be read as a caution to re-tune it whenever the closed form changes, not as a claim that every reported variant used its individually best setting.

\textbf{3. Tree-ensemble comparison and hybrid scope.} The Gradient-Boosted Trees hyperparameter search (Section~5.4, 36 combinations) is a modest grid, and Random Forest was left untuned entirely for comparison; that a 36-point search barely moved Gradient-Boosted Trees' numbers suggests little is left nearby in hyperparameter space, but a larger search, or tuning Random Forest too, might still close a little more of the gap. Both physics-plus-boosting hybrids build on only one PGML variant (\emph{2-channel}) with hyperparameters reused from the standalone tree baseline; we did not try either hybrid on the other three PGML variants or the symmetric hybrid (a closed-form correction on top of a tree ensemble's prediction), any of which might change how much of the gap closes.

\textbf{4. Ensemble strategy and pool-size scope.} The three combination strategies compared (unweighted average, Ridge stacking at a single fixed $\alpha=1$, inverse-variance weighting) do not exhaust the space of ways to combine models, so none of the reported results is necessarily the best any strategy could do; the stacked variant's instability reflects the same small-validation-split limitation above (12 curves per seed to fit 8--10 regression weights), and a per-seed-tuned regularization strength or a shrinkage estimator between the fully free stack and the one-parameter-per-model scheme might close some of the remaining gap to it. The ten-model pool (Section~5.6) was checked as one enlarged set rather than via individually added hybrids or a hand-picked smaller subset, and the inverse-variance weighting scheme assumes base-model errors are uncorrelated -- an assumption its own correlated members (four related PGML variants, two tree ensembles) violate to some degree -- without comparison against other low-parameter alternatives (e.g.\ rank-based weighting); its advantage over stacking should be read as evidence that \emph{a} low-parameter scheme beats a high-parameter one here, not that this specific scheme is the best possible one.

\textbf{5. Everything above is computational.} All evaluation here is computational, on synthetic data generated by a model in the same functional family the paper's own closed-form revisions are then credited with matching; nothing in this paper should be read as a claim about which modeling strategy is preferable on real, measured semiconductor conductivity data. Section~\ref{sec:conclusion} discusses what would settle this.

\section{Conclusion}
\label{sec:conclusion}

\textbf{What we found.} We set out to test, rather than assume, whether routing a network's predictions through a closed-form transport layer improves on a black-box network and a physics-regularized network with no such layer. Across seven successive revisions on a synthetic benchmark parameterized by 24 real semiconductors' descriptors, the answer depended entirely on the closed form's functional richness: a reduced closed form was the worst model on every metric, including the physical-consistency violation rate it was meant to improve, while a richer closed form matching the benchmark generator's own functional family resolved consistency completely and closed most of the accuracy gap. Neither revision closed the remaining gap to Gradient-Boosted Trees, which also survived direct hyperparameter tuning; a physics-plus-boosting hybrid narrowed it only partly, while reaching the best extrapolation robustness of any single model on this synthetic benchmark. Combining every model tried, weighted by validation accuracy, reached the lowest mean error and best mean extrapolation robustness of anything tested on this synthetic benchmark, and neither widening the combination's candidate pool nor stacking instead of a simple weighting scheme improved on it further.

\textbf{What this does not show.} Three qualifications bound every result above. First, the richer closed form's success is closer to a demonstration of correct model specification than of general architectural benefit: it was deliberately built in the same functional family as the benchmark's own generator, so we cannot separate "the architecture helps" from "the architecture was handed the right functional form to fit." Second, the final ensemble's numerical edge over tuned Gradient-Boosted Trees is not statistically significant in a paired test across the 5 seeds on any metric ($p > 0.1$), nor in the exploratory material-level bootstrap of Section~4.2; both tests have limited power at this benchmark's scale (5 seeds, 24 materials), and the result should be read as suggestive, not established. Third, and independently of significance, that ensemble includes Gradient-Boosted Trees itself as a member, so even a significant edge would not be evidence that the physics-guided architecture outperforms it unaided -- no PGML variant, in isolation, closed that gap at any point in this study. Whether architectural physics integration can win outright on this benchmark, rather than by being folded into a mixture with the baseline it is compared against, remains exactly as open as it was after the third experiment.

\textbf{What would settle it.} All data here are synthetic, generated by the same reduced transport model the richer closed form was built to match, and the benchmark spans only 24 materials with 3--5 held out per seed. The next test this paper's design points to directly is the same seven-experiment comparison run on real, measured $\sigma(T)$ data, with more materials, more seeds, and the same paired significance testing applied throughout -- which would show whether either qualification above still holds once the closed form's advantage is not partly a byproduct of matching its own data-generating process. We release the benchmark generator, every model and baseline implementation, the training and significance-testing code, and per-seed results so that this comparison can be run.

\section*{Code and Data Availability}
The full synthetic-benchmark generator, all closed-form model implementations (the current \texttt{code/models.py} implements the richer base closed form with the explicit second channel, Table~\ref{tab:richer}; see \texttt{code/README.md} for how the more reduced closed form behind Table~\ref{tab:main} differed) and every baseline, and the training and multi-seed evaluation scripts used to produce every number in this paper are released alongside the manuscript in \texttt{code/}. Running \texttt{python run\_multi\_seed.py} against the current \texttt{code/models.py} reproduces Tables~\ref{tab:richer}, \ref{tab:ablation}, \ref{tab:hybrid}, \ref{tab:ensemble}, \ref{tab:ensemble2}, and \ref{tab:ensemble3}. Gradient-Boosted Trees' hyperparameter search, both physics-plus-boosting hybrids of Section~5.4, and all six ensemble variants of Sections~5.5--5.7 are all computed in that same script. \texttt{code/results\_seed*.json} and \texttt{code/results\_multiseed.json} record the exact run each table is drawn from, including the per-seed stacked-ensemble weights and validation-RMSE weights; \texttt{code/significance\_tests.py} recomputes the paired $t$-test and Wilcoxon signed-rank results reported in Sections~5.5--5.7 and the Conclusion directly from those per-seed files, and \texttt{code/bootstrap\_tests.py} recomputes the material-level bootstrap confidence intervals and $p$-values reported alongside them (Section~4.2) from the same files' per-curve RMSE breakdown. The electronic-structure descriptors are representative literature values in the style of public repositories including the Materials Project \citep{jain2013materialsproject} and AFLOW \citep{curtarolo2012aflow}, tabulated directly in \texttt{code/descriptors.py} (Section~3.2 documents which descriptors are exact computations versus literature-typical values or engineered proxies). No experimental or proprietary data were used. The repository is version-controlled with a full commit history, which identifies the exact state of the code and text corresponding to this manuscript; a permanent archival identifier (e.g.\ a Zenodo DOI tied to a tagged release) is not yet assigned. Each of the 5 random seeds used throughout is fixed and given explicitly in \texttt{code/run\_multi\_seed.py} so that every reported run is exactly reproducible from a clean checkout.

\section*{Author Contributions (CRediT)}
Conceptualization, Methodology, Software, Formal Analysis, Writing -- Original Draft, Writing -- Review \& Editing: M.H.K.

\section*{Conflict of Interest Statement}
The author declares no conflict of interest.

\section*{Funding Statement}
This research received no external funding.

\bibliographystyle{unsrtnat}
\bibliography{references}

\end{document}
